\documentclass[10pt,a4paper]{amsart}
\usepackage{amsmath,amssymb,amsthm,graphicx}
\usepackage{verbatim}
\newenvironment{reference}{\comment}{\endcomment}
%\newenvironment{reference}{}{}
\newenvironment{slogan}{\comment}{\endcomment}
\newenvironment{history}{\comment}{\endcomment}

% For commutative diagrams we use Xy-pic
\usepackage[all]{xy}

% We use 2cell for 2-commutative diagrams.
\xyoption{2cell}
% Theorem environments.
%
\theoremstyle{plain}
\newtheorem{theorem}[subsection]{Theorem}
\newtheorem{proposition}[subsection]{Proposition}
\newtheorem{lemma}[subsection]{Lemma}

\theoremstyle{definition}
\newtheorem{definition}[subsection]{Definition}
\newtheorem{example}[subsection]{Example}
\newtheorem{exercise}[subsection]{Exercise}
\newtheorem{situation}[subsection]{Situation}

\theoremstyle{remark}
\newtheorem{remark}[subsection]{Remark}
\newtheorem{remarks}[subsection]{Remarks}

\numberwithin{equation}{subsection}

% Macros
%
\def\lim{\mathop{\mathrm{lim}}\nolimits}
\def\colim{\mathop{\mathrm{colim}}\nolimits}
\def\Spec{\mathop{\mathrm{Spec}}}
\def\Hom{\mathop{\mathrm{Hom}}\nolimits}
\def\Ext{\mathop{\mathrm{Ext}}\nolimits}
\def\SheafHom{\mathop{\mathcal{H}\!\mathit{om}}\nolimits}
\def\SheafExt{\mathop{\mathcal{E}\!\mathit{xt}}\nolimits}
\def\Sch{\mathit{Sch}}
\def\Mor{\mathop{\mathrm{Mor}}\nolimits}
\def\Ob{\mathop{\mathrm{Ob}}\nolimits}
\def\Sh{\mathop{\mathit{Sh}}\nolimits}
\def\NL{\mathop{N\!L}\nolimits}
\def\CH{\mathop{\mathrm{CH}}\nolimits}
\def\proetale{{pro\text{-}\acute{e}tale}}
\def\etale{{\acute{e}tale}}
\def\QCoh{\mathit{QCoh}}
\def\Ker{\mathop{\mathrm{Ker}}}
\def\Im{\mathop{\mathrm{Im}}}
\def\Coker{\mathop{\mathrm{Coker}}}
\def\Coim{\mathop{\mathrm{Coim}}}

% Boxtimes
%
\DeclareMathSymbol{\boxtimes}{\mathbin}{AMSa}{"02}

%
% Macros for moduli stacks/spaces
%
\def\QCohstack{\mathcal{QC}\!\mathit{oh}}
\def\Cohstack{\mathcal{C}\!\mathit{oh}}
\def\Spacesstack{\mathcal{S}\!\mathit{paces}}
\def\Quotfunctor{\mathrm{Quot}}
\def\Hilbfunctor{\mathrm{Hilb}}
\def\Curvesstack{\mathcal{C}\!\mathit{urves}}
\def\Polarizedstack{\mathcal{P}\!\mathit{olarized}}
\def\Complexesstack{\mathcal{C}\!\mathit{omplexes}}
% \Pic is the operator that assigns to X its picard group, usage \Pic(X)
% \Picardstack_{X/B} denotes the Picard stack of X over B
% \Picardfunctor_{X/B} denotes the Picard functor of X over B
\def\Pic{\mathop{\mathrm{Pic}}\nolimits}
\def\Picardstack{\mathcal{P}\!\mathit{ic}}
\def\Picardfunctor{\mathrm{Pic}}
\def\Deformationcategory{\mathcal{D}\!\mathit{ef}}
\usepackage{cleveref}
\crefname{originalreference}{original reference}{original references}
\Crefname{originalreference}{Original reference}{Original references}
\makeatletter
\newlabel{lemma-compute-what-it-should-be-derived}{{0DIC}{source}{}{sourcepointer-1}{}}
\newlabel{lemma-compute-what-it-should-be-derived@cref}{{[originalreference][0][]0DIC}{[1][0][]source}{}{}{}}
\makeatother
\crefname{originallemma}{original lemma}{original lemmas}
\Crefname{originallemma}{Original lemma}{Original lemmas}
\makeatletter
\newlabel{more-algebra-lemma-Rlim-pseudo-coherent-gives-pseudo-coherent}{{0CQF}{source}{}{sourcepointer-2}{}}
\newlabel{more-algebra-lemma-Rlim-pseudo-coherent-gives-pseudo-coherent@cref}{{[originallemma][0][]0CQF}{[1][0][]source}{}{}{}}
\makeatother
\crefname{originalreference}{original reference}{original references}
\Crefname{originalreference}{Original reference}{Original references}
\makeatletter
\newlabel{situation-existence-derived}{{0DIB}{source}{}{sourcepointer-3}{}}
\newlabel{situation-existence-derived@cref}{{[originalreference][0][]0DIB}{[1][0][]source}{}{}{}}
\makeatother
\crefname{originallemma}{original lemma}{original lemmas}
\Crefname{originallemma}{Original lemma}{Original lemmas}
\makeatletter
\newlabel{spaces-perfect-lemma-flat-proper-pseudo-coherent-direct-image-general}{{0CTN}{source}{}{sourcepointer-4}{}}
\newlabel{spaces-perfect-lemma-flat-proper-pseudo-coherent-direct-image-general@cref}{{[originallemma][0][]0CTN}{[1][0][]source}{}{}{}}
\makeatother
\title{Derived global sections against a perfect object are pseudo-coherent and commute with base change along the derived descent tower}
\author{The Stacks Project Authors}
\date{stacks-project:spaces-flat.tex:local-snapshot-20261012; GNU FDL 1.2 or later}
\begin{document}
\maketitle
\noindent\textit{Editorial note.} All delivered mathematical text is the author's own, copied from the pinned Stacks Project chapter file spaces-flat.tex; no statement or proof step was rewritten, reordered or removed.\par\medskip
\noindent\textit{Editorial note.} The counted proof is the author's complete argument: a stage-wise RGamma computation with base change compatible pseudo-coherent transitions, the derived-limit commutation of RGamma(X, -), and the algebra limit lemma for pseudo-coherent objects\par\medskip
\noindent\textit{Editorial note.} The label lemma-compute-against-perfect-derived names the counted statement, and the proof environment immediately following it in the source is delivered verbatim as its proof.\par\medskip
\noindent\textit{Editorial note.} The counted statement is carried by the author's own lemma environment in the chapter source.\par\medskip
\noindent\textit{Allowed standard prerequisites (editorial).}\begin{itemize}
\item The derived descent tower situation with its system of pseudo-coherent objects K\_n and the limit object K defined as the derived limit in the quasi-coherent derived category
\item The derived categories of spaces lemma that the direct image computing the stages of a pseudo-coherent system is pseudo-coherent with formation commuting with base change
\item The More on Algebra lemma that a derived limit of pseudo-coherent objects with compatible base change data is pseudo-coherent and recovers the stages on base change
\item The exact functor RGamma on the quasi-coherent derived category commuting with products and hence with derived limits
\end{itemize}\medskip
\ifdefined\hypertarget\hypertarget{sourcepointer-1}{}\fi
\noindent\textit{Source reference.} Original reference 0DIC: sources/stacks-spaces-flat/pinned-original-source.tex, source lines 2695--2695. Prerequisite pointer invoked at proof line(s) 2755; it names a standard fact of the same chapter that the counted argument uses, and every step of the counted proof built on it is retained verbatim in the delivered body.\par\medskip
\ifdefined\hypertarget\hypertarget{sourcepointer-2}{}\fi
\noindent\textit{Source reference.} Original lemma 0CQF: sources/stacks-chow/tags.txt, source lines 15960--15960. Prerequisite pointer invoked at proof line(s) 2777; it names a standard fact of a sibling chapter that the counted argument uses, and every step of the counted proof built on it is retained verbatim in the delivered body.\par\medskip
\ifdefined\hypertarget\hypertarget{sourcepointer-3}{}\fi
\noindent\textit{Source reference.} Original reference 0DIB: sources/stacks-spaces-flat/pinned-original-source.tex, source lines 2668--2668. Prerequisite pointer invoked at proof line(s) 2754; it names a standard fact of the same chapter that the counted argument uses, and every step of the counted proof built on it is retained verbatim in the delivered body.\par\medskip
\ifdefined\hypertarget\hypertarget{sourcepointer-4}{}\fi
\noindent\textit{Source reference.} Original lemma 0CTN: sources/stacks-chow/tags.txt, source lines 16073--16073. Prerequisite pointer invoked at proof line(s) 2771; it names a standard fact of a sibling chapter that the counted argument uses, and every step of the counted proof built on it is retained verbatim in the delivered body.\par\medskip
\begin{lemma}
\label{lemma-compute-against-perfect-derived}
In Situation \ref{situation-existence-derived} let $K$ be as in
Lemma \ref{lemma-compute-what-it-should-be-derived}. For any perfect
object $E$ of $D(\mathcal{O}_X)$ the cohomology
$$
M = R\Gamma(X, K \otimes^\mathbf{L} E)
$$
is a pseudo-coherent object of $D(A)$ and there is a canonical isomorphism
$$
R\Gamma(X_n, K_n \otimes^\mathbf{L} E|_{X_n}) = M \otimes_A^\mathbf{L} A_n
$$
in $D(A_n)$. Here $E|_{X_n}$ denotes the derived pullback of $E$ to $X_n$.
\end{lemma}

\begin{proof}
Write $E_n = E|_{X_n}$ and
$M_n = R\Gamma(X_n, K_n \otimes^\mathbf{L} E|_{X_n})$.
By Derived Categories of Spaces, Lemma
\ref{spaces-perfect-lemma-flat-proper-pseudo-coherent-direct-image-general}
we see that $M_n$ is a pseudo-coherent object of $D(A_n)$
whose formation commutes with base change. Thus the maps
$M_n \otimes_{A_n}^\mathbf{L} A_{n - 1} \to M_{n - 1}$
coming from $\varphi_n$ are isomorphisms. By
More on Algebra, Lemma
\ref{more-algebra-lemma-Rlim-pseudo-coherent-gives-pseudo-coherent}
we find that $R\lim M_n$ is pseudo-coherent and
that its base change back to $A_n$ recovers $M_n$.
On the other hand, the exact functor
$R\Gamma(X, -) : D_\QCoh(\mathcal{O}_X) \to D(A)$
of triangulated categories commutes with products
and hence with derived limits, whence
$$
R\Gamma(X, E \otimes^\mathbf{L} K) =
R\lim R\Gamma(X,  E \otimes^\mathbf{L} K_n) =
R\lim R\Gamma(X_n, E_n \otimes^\mathbf{L} K_n) =
R\lim M_n
$$
as desired.
\end{proof}

\end{document}
